% LaTeX Template for AI Problem Analysis Output (v2.0)
% AMC 12B 2023 Problem 15 Strategic Analysis
% Structure your analysis results using this template with colorful boxes
% IMPORTANT: Use LaTeX commands for all mathematical symbols, NOT unicode characters

\documentclass[12pt]{article}
\usepackage[utf8]{inputenc}
\usepackage{amsmath, amssymb, amsthm}
\usepackage{geometry}
\usepackage{xcolor}
\usepackage[most]{tcolorbox}
\usepackage{enumitem}
\usepackage{fancyhdr}

% Page setup
\geometry{margin=1in}
\setlength{\headheight}{14.5pt}
\pagestyle{fancy}
\fancyhf{}
\rhead{AMC12 Problem Analysis}
\lhead{\today}

% Color definitions
\definecolor{problemconfig}{RGB}{230, 242, 255}    % Light blue
\definecolor{strategic}{RGB}{255, 230, 242}       % Light pink
\definecolor{tactical}{RGB}{242, 255, 230}        % Light green
\definecolor{techniques}{RGB}{255, 242, 230}      % Light yellow
\definecolor{verification}{RGB}{255, 230, 230}    % Light red
\definecolor{solution}{RGB}{230, 255, 255}        % Light cyan
\definecolor{competition}{RGB}{242, 242, 255}     % Light purple
\definecolor{proofwriting}{RGB}{240, 230, 255}    % Light lavender
\definecolor{gold}{RGB}{218, 165, 32}

% Box styles
\newtcolorbox{problemconfigbox}{
    colback=problemconfig,
    colframe=blue,
    title=Problem Configuration Analysis,
    fonttitle=\bfseries,
    arc=3pt,
    boxrule=1pt,
    breakable
}

\newtcolorbox{strategicbox}{
    colback=strategic,
    colframe=purple,
    title=Strategic Framework Application,
    fonttitle=\bfseries,
    arc=3pt,
    boxrule=1pt,
    breakable
}

\newtcolorbox{tacticalbox}{
    colback=tactical,
    colframe=green,
    title=Tactical Methodology Selection,
    fonttitle=\bfseries,
    arc=3pt,
    boxrule=1pt,
    breakable
}

\newtcolorbox{techniquesbox}{
    colback=techniques,
    colframe=orange,
    title=Conversion Techniques Application,
    fonttitle=\bfseries,
    arc=3pt,
    boxrule=1pt,
    breakable
}

\newtcolorbox{verificationbox}{
    colback=verification,
    colframe=red,
    title=Verification Strategy Design,
    fonttitle=\bfseries,
    arc=3pt,
    boxrule=1pt,
    breakable
}

\newtcolorbox{solutionbox}{
    colback=solution,
    colframe=cyan,
    title=Solution Roadmap,
    fonttitle=\bfseries,
    arc=3pt,
    boxrule=1pt,
    breakable
}

\newtcolorbox{competitionbox}{
    colback=competition,
    colframe=purple,
    title=Competition Strategy Integration,
    fonttitle=\bfseries,
    arc=3pt,
    boxrule=1pt,
    breakable
}

\newtcolorbox{keyinsightbox}{
    colback=yellow!20,
    colframe=gold,
    title=Key Insight,
    fonttitle=\bfseries,
    arc=3pt,
    boxrule=2pt,
    leftrule=5pt,
    breakable
}

\newtcolorbox{warningbox}{
    colback=red!10,
    colframe=red!50,
    title=Warning: Common Pitfall,
    fonttitle=\bfseries,
    arc=3pt,
    boxrule=2pt,
    leftrule=5pt,
    breakable
}

\newtcolorbox{tipbox}{
    colback=green!10,
    colframe=green!50,
    title=Pro Tip,
    fonttitle=\bfseries,
    arc=3pt,
    boxrule=2pt,
    leftrule=5pt,
    breakable
}

\begin{document}

\title{AMC12 Strategic Problem Analysis\\2023 AMC 12B Problem 15}
\author{Competition Math Framework v2.1}
\date{\today}
\maketitle

\section*{Problem Statement}

Suppose $a$, $b$, and $c$ are positive integers such that
$$\frac{a}{14}+\frac{b}{15}=\frac{c}{210}.$$

Which of the following statements are necessarily true?

\textbf{I.} If $\gcd(a,14)=1$ or $\gcd(b,15)=1$ or both, then $\gcd(c,210)=1$.

\textbf{II.} If $\gcd(c,210)=1$, then $\gcd(a,14)=1$ or $\gcd(b,15)=1$ or both.

\textbf{III.} $\gcd(c,210)=1$ if and only if $\gcd(a,14)=\gcd(b,15)=1$.

$$\textbf{(A)}~\text{I, II, and III}\qquad\textbf{(B)}~\text{I only}\qquad\textbf{(C)}~\text{I and II only}\qquad\textbf{(D)}~\text{III only}\qquad\textbf{(E)}~\text{II and III only}$$

\begin{problemconfigbox}
\subsection*{A. Problem Classification}
\begin{itemize}[leftmargin=*]
    \item \textbf{Problem Type}: To Prove/Disprove (determine which logical statements are true)
    \item \textbf{Competition Context}: AMC12 2023B Problem 15
    \item \textbf{Difficulty Band}: Mid-to-Late transition (P15/25) - requires number theory reasoning and logical analysis
    \item \textbf{Mathematical Domain}: Number Theory (GCD properties) / Logic (if-then, if-and-only-if statements)
    \item \textbf{Estimated Time}: 5-7 minutes for careful analysis or strategic testing
\end{itemize}

\subsection*{B. Problem Structure Identification}
\begin{itemize}[leftmargin=*]
    \item \textbf{Given Information}: 
    \begin{itemize}
        \item Equation: $\frac{a}{14}+\frac{b}{15}=\frac{c}{210}$
        \item Equivalent form: $15a + 14b = c$ (multiply by $\text{lcm}(14,15) = 210$)
        \item $a, b, c$ are positive integers
        \item Three logical statements about GCD relationships
        \item Five answer choices combining the statements
    \end{itemize}
    \item \textbf{Constraints}: 
    \begin{itemize}
        \item $a, b, c > 0$ (positive integers)
        \item Linear Diophantine equation: $15a + 14b = c$
        \item Note: $14 = 2 \cdot 7$, $15 = 3 \cdot 5$, $210 = 2 \cdot 3 \cdot 5 \cdot 7 = 14 \cdot 15$
        \item $\gcd(14, 15) = 1$ (coprime)
    \end{itemize}
    \item \textbf{Objective}: Determine which statements are necessarily true (always true for all valid $a, b, c$)
    \item \textbf{Hidden Assumptions}: 
    \begin{itemize}
        \item Need to distinguish between ``or'' (at least one) and ``and'' (both)
        \item ``If and only if'' means bidirectional: $P \Leftrightarrow Q$ means $P \Rightarrow Q$ AND $Q \Rightarrow P$
        \item Can disprove statements with counterexamples
        \item Must prove true statements rigorously or use GCD properties
    \end{itemize}
    \item \textbf{Answer Choice Leverage}: 
    \begin{itemize}
        \item Five choices with different combinations
        \item If we can quickly disprove Statement I, we eliminate (A), (B), (C)
        \item Statement III appears in both (D) and (E), so it's likely true
        \item The difference between (D) and (E) is Statement II
    \end{itemize}
\end{itemize}

\subsection*{C. Strategic Orientation}
\begin{itemize}[leftmargin=*]
    \item \textbf{Hypothesis}: Statement I is likely false (easiest to find counterexample), Statements II and III are likely true
    \item \textbf{Conclusion}: Determine which combination of I, II, III is correct
    \item \textbf{Notation}: 
    \begin{itemize}
        \item Let $d_a = \gcd(a, 14)$, $d_b = \gcd(b, 15)$, $d_c = \gcd(c, 210)$
        \item Key equation: $15a + 14b = c$
        \item Prime factorizations: $14 = 2 \cdot 7$, $15 = 3 \cdot 5$, $210 = 2 \cdot 3 \cdot 5 \cdot 7$
    \end{itemize}
    \item \textbf{Intuitive Approach}: 
    \begin{enumerate}
        \item Rewrite equation in cleaner form: $15a + 14b = c$
        \item Test Statement I with small counterexamples
        \item Use GCD properties and modular arithmetic for Statements II and III
        \item Leverage answer choices to narrow possibilities
    \end{enumerate}
    \item \textbf{Cross-Domain Potential}: This combines number theory (GCD), logic (conditional statements), and algebra (Diophantine equations).
\end{itemize}
\end{problemconfigbox}

\begin{strategicbox}
\subsection*{A. Psychological Strategies}
\begin{itemize}[leftmargin=*]
    \item \textbf{Mental Toughness}: Multi-part logic problems require careful reading. Don't rush - understand what each statement claims.
    \item \textbf{Creativity Cultivation}: Finding counterexamples requires creativity. Try simple values first ($a=1, b=1$, etc.).
    \item \textbf{Iterative Refinement}: If a statement seems false, try a few counterexamples. If they all fail, maybe the statement is true!
\end{itemize}

\subsection*{B. Investigation Strategies}
\begin{itemize}[leftmargin=*]
    \item \textbf{Startup Orientation}: Immediately convert the fraction equation to $15a + 14b = c$ for clarity.
    \item \textbf{Get Your Hands Dirty}: Test Statement I with small values like $a=1, b=1$ or $a=2, b=2$.
    \item \textbf{Penultimate Step}: Once you narrow to two answer choices, focus on the distinguishing statement.
    \item \textbf{Make It Easier}: 
    \begin{itemize}
        \item Use answer choice structure: If I is false, eliminate (A), (B), (C)
        \item Notice III appears in both remaining choices, so it must be true
        \item Only need to determine if II is true
    \end{itemize}
    \item \textbf{Wishful Thinking}: Wish for a simple counterexample to I - often the simplest values work!
    \item \textbf{Conjecture via Small Cases}: Test with $a=1, 2, 3$ and $b=1, 2, 3$ to see patterns.
\end{itemize}

\subsection*{C. Argument Construction}
\begin{itemize}[leftmargin=*]
    \item \textbf{Proof Method}: 
    \begin{itemize}
        \item For false statements: Find counterexamples
        \item For true statements: Use GCD properties and modular arithmetic
        \item For ``if and only if'': Prove both directions
    \end{itemize}
    \item \textbf{Key Lemmas}: 
    \begin{enumerate}
        \item $15a + 14b = c$
        \item GCD property: $\gcd(ma + nb, m) = \gcd(nb, m) = \gcd(b, m) \cdot \gcd(n, m/\gcd(b,m))$
        \item Simplified: $\gcd(c, 14) = \gcd(15a + 14b, 14) = \gcd(15a, 14) = \gcd(a, 14)$ (since $\gcd(15,14)=1$)
        \item Similarly: $\gcd(c, 15) = \gcd(b, 15)$
        \item Therefore: $\gcd(c, 210) = \gcd(c, 14) \cdot \gcd(c, 15) = \gcd(a, 14) \cdot \gcd(b, 15)$ (since $\gcd(14,15)=1$)
    \end{enumerate}
    \item \textbf{Logical Structure}: Test statements systematically, use answer choices to guide strategy
\end{itemize}
\end{strategicbox}

\begin{tacticalbox}
\subsection*{A. Core Mathematical Tactics}
\begin{itemize}[leftmargin=*]
    \item \textbf{Primary Tactics}: 
    \begin{itemize}
        \item \emph{Counterexample Testing}: Find specific values that violate a statement
        \item \emph{GCD Properties}: Use properties of greatest common divisor
        \item \emph{Modular Arithmetic}: Take equation modulo prime factors
        \item \emph{Logical Analysis}: Understand ``or'' vs ``and'', ``if-then'' vs ``if-and-only-if''
    \end{itemize}
    \item \textbf{Domain-Specific Tactics (Number Theory)}: 
    \begin{itemize}
        \item GCD multiplication rule: If $\gcd(m,n)=1$, then $\gcd(a, mn) = \gcd(a,m) \cdot \gcd(a,n)$
        \item Linear combination: $\gcd(ma+nb, n) = \gcd(ma, n)$ when appropriate
        \item Prime factorization: $14 = 2 \cdot 7$, $15 = 3 \cdot 5$, $210 = 2 \cdot 3 \cdot 5 \cdot 7$
        \item Coprimality: $\gcd(14,15) = 1$ is crucial
    \end{itemize}
    \item \textbf{Domain-Specific Tactics (Logic)}: 
    \begin{itemize}
        \item ``Or'' means at least one (inclusive or)
        \item ``If and only if'' requires proving both directions
        \item Contrapositive: ``If P then Q'' is equivalent to ``If not Q then not P''
    \end{itemize}
\end{itemize}
\end{tacticalbox}

\begin{techniquesbox}
\subsection*{A. High-Probability Techniques}
\begin{itemize}[leftmargin=*]
    \item \textbf{Primary Technique}: \textbf{Answer Choice Testing} + \textbf{GCD Properties}
    \item \textbf{Recognition Cues}: 
    \begin{itemize}
        \item Multiple logical statements to evaluate
        \item Answer choices combine statements in different ways
        \item Counterexamples can eliminate statements
        \item Small numbers ($14, 15, 210$) suggest testing is feasible
    \end{itemize}
    \item \textbf{Application - Strategic Approach}: 
    \begin{enumerate}
        \item Convert equation: $\frac{a}{14}+\frac{b}{15}=\frac{c}{210} \Rightarrow 15a + 14b = c$
        \item Test Statement I with simple counterexample: Try $a=2, b=2$
        \begin{itemize}
            \item $c = 15(2) + 14(2) = 30 + 28 = 58$
            \item $\gcd(2, 14) = 2 \neq 1$, $\gcd(2, 15) = 1$ $\checkmark$
            \item So I's condition is satisfied (at least one GCD is 1)
            \item But $\gcd(58, 210) = \gcd(58, 210) = 2 \neq 1$
            \item Statement I is FALSE!
        \end{itemize}
        \item Eliminate answer choices (A), (B), (C)
        \item Remaining choices: (D) III only, or (E) II and III only
        \item Since III appears in both, Statement III must be TRUE
        \item Determine if Statement II is true to choose between (D) and (E)
    \end{enumerate}
\end{itemize}

\subsection*{B. Alternative Techniques}
\begin{itemize}[leftmargin=*]
    \item \textbf{Rigorous GCD Analysis}: 
    \begin{itemize}
        \item Use GCD property: $\gcd(c, 14) = \gcd(15a+14b, 14) = \gcd(15a, 14) = \gcd(a, 14)$
        \item Similarly: $\gcd(c, 15) = \gcd(b, 15)$
        \item Since $\gcd(14,15)=1$: $\gcd(c, 210) = \gcd(a,14) \cdot \gcd(b,15)$
        \item This directly proves Statement III!
    \end{itemize}
    \item \textbf{Modular Arithmetic Approach}: 
    \begin{itemize}
        \item For Statement III's ``if'' direction: Assume $\gcd(a,14)=\gcd(b,15)=1$ but $\gcd(c,210) \neq 1$
        \item Then $c$ has a prime factor $p \in \{2,3,5,7\}$
        \item Taking modulo $p$ on $15a+14b=c$: contradiction arises
    \end{itemize}
    \item \textbf{Technique Integration}: Answer choice elimination is fastest, but GCD analysis provides rigorous proof.
\end{itemize}
\end{techniquesbox}

\begin{verificationbox}
\subsection*{A. Verification Level Selection}
\begin{itemize}[leftmargin=*]
    \item \textbf{Chosen Level}: Level 4-5 (Standard to Thorough, 1-2 minutes)
    \item \textbf{Justification}: Logic problem with multiple statements. Need to verify counterexample works and check remaining statements.
    \item \textbf{Time Budget}: 1-2 minutes for verification
\end{itemize}

\subsection*{B. Verification Checklist}
\begin{itemize}[leftmargin=*]
    \item \textbf{Counterexample Validity}: 
    \begin{itemize}
        \item Does counterexample satisfy the original equation?
        \item Does it satisfy the condition of the statement?
        \item Does it violate the conclusion of the statement?
    \end{itemize}
    \item \textbf{Logical Rigor}: 
    \begin{itemize}
        \item For ``or'' statements: Is at least one condition satisfied?
        \item For ``and'' statements: Are both conditions satisfied?
        \item For ``if and only if'': Are both directions proven?
    \end{itemize}
    \item \textbf{Answer Choice Check}: Does selection match one of the five options?
\end{itemize}

\subsection*{C. Detailed Verification}

\textbf{Verify Statement I is False}:

Counterexample: $a=2, b=2$
\begin{align*}
    c &= 15(2) + 14(2) = 58\\
    \gcd(2, 14) &= 2 \neq 1\\
    \gcd(2, 15) &= 1 \checkmark \text{ (at least one GCD is 1)}\\
    \gcd(58, 210) &= \gcd(58, 210)
\end{align*}

Find $\gcd(58, 210)$:
\begin{align*}
    58 &= 2 \cdot 29\\
    210 &= 2 \cdot 3 \cdot 5 \cdot 7\\
    \gcd(58, 210) &= 2 \neq 1
\end{align*}

Statement I claims: condition satisfied $\Rightarrow$ $\gcd(c,210)=1$

But we have: condition satisfied, yet $\gcd(c,210)=2 \neq 1$

Therefore Statement I is FALSE $\checkmark$
\end{verificationbox}

\begin{solutionbox}
\subsection*{A. Step-by-Step Solution (Strategic Approach)}

\begin{enumerate}[leftmargin=*]
    \item \textbf{Convert equation to cleaner form (30 seconds)}
    
    Given: $\frac{a}{14}+\frac{b}{15}=\frac{c}{210}$
    
    Multiply both sides by $210 = \text{lcm}(14, 15)$:
    $$210 \cdot \frac{a}{14} + 210 \cdot \frac{b}{15} = 210 \cdot \frac{c}{210}$$
    $$15a + 14b = c$$
    
    This is our working equation.
    
    \item \textbf{Test Statement I with counterexample (60 seconds)}
    
    \textbf{Statement I}: If $\gcd(a,14)=1$ or $\gcd(b,15)=1$ or both, then $\gcd(c,210)=1$.
    
    Try $a=2, b=2$:
    \begin{align*}
        c &= 15(2) + 14(2) = 30 + 28 = 58\\
        \gcd(2, 14) &= \gcd(2, 2 \cdot 7) = 2 \neq 1\\
        \gcd(2, 15) &= \gcd(2, 3 \cdot 5) = 1 \checkmark
    \end{align*}
    
    Condition is satisfied: $\gcd(b,15)=1$ $\checkmark$
    
    Check conclusion:
    \begin{align*}
        \gcd(58, 210) &= \gcd(2 \cdot 29, 2 \cdot 3 \cdot 5 \cdot 7) = 2 \neq 1
    \end{align*}
    
    Statement I is \textbf{FALSE}! (Counterexample found)
    
    \item \textbf{Eliminate answer choices (10 seconds)}
    
    Statement I is false, so eliminate:
    \begin{itemize}
        \item (A) I, II, and III - \textbf{Eliminated}
        \item (B) I only - \textbf{Eliminated}
        \item (C) I and II only - \textbf{Eliminated}
    \end{itemize}
    
    Remaining choices:
    \begin{itemize}
        \item (D) III only
        \item (E) II and III only
    \end{itemize}
    
    \item \textbf{Observe that Statement III must be true (15 seconds)}
    
    Since Statement III appears in both remaining answer choices (D) and (E), it \emph{must} be true.
    
    The only difference is whether Statement II is true.
    
    \item \textbf{Determine Statement II (2-3 minutes)}
    
    \textbf{Statement II}: If $\gcd(c,210)=1$, then $\gcd(a,14)=1$ or $\gcd(b,15)=1$ or both.
    
    We can prove this using the contrapositive or by using GCD properties.
    
    \textbf{Method 1: Contrapositive}
    
    Contrapositive: If $\gcd(a,14) \neq 1$ AND $\gcd(b,15) \neq 1$, then $\gcd(c,210) \neq 1$.
    
    Assume $\gcd(a,14) > 1$ and $\gcd(b,15) > 1$.
    
    Then:
    \begin{itemize}
        \item $\gcd(a,14) \geq 2$ (or 7), so $a$ shares a factor with 14
        \item $\gcd(b,15) \geq 3$ (or 5), so $b$ shares a factor with 15
    \end{itemize}
    
    Case: Suppose $2|a$ (since $2|14$).
    
    From $c = 15a + 14b$:
    \begin{align*}
        c &\equiv 15a + 14b \pmod{2}\\
        &\equiv a + 0 \pmod{2} \quad \text{(since $15 \equiv 1$, $14 \equiv 0 \pmod{2}$)}\\
        &\equiv 0 \pmod{2} \quad \text{(since $2|a$)}
    \end{align*}
    
    So $2|c$, thus $\gcd(c,210) \geq 2 > 1$ $\checkmark$
    
    Similar arguments work for other prime factors.
    
    Therefore Statement II is \textbf{TRUE}.
    
    \item \textbf{Verify Statement III for completeness (optional, 2-3 minutes)}
    
    \textbf{Statement III}: $\gcd(c,210)=1$ if and only if $\gcd(a,14)=\gcd(b,15)=1$.
    
    This requires proving both directions.
    
    \textbf{``If'' direction ($\Rightarrow$)}: 
    
    If $\gcd(a,14)=1$ and $\gcd(b,15)=1$, then $\gcd(c,210)=1$.
    
    Suppose not: $\gcd(c,210) > 1$. Then $c$ has a prime factor $p \in \{2,3,5,7\}$.
    
    \begin{itemize}
        \item If $p \in \{2,7\}$ (factors of 14): Then $p|c$ implies $p|15a$ (from $c=15a+14b$). Since $p \nmid 15$, we have $p|a$. But $\gcd(a,14)=1$ means $p \nmid a$. Contradiction.
        \item If $p \in \{3,5\}$ (factors of 15): Then $p|c$ implies $p|14b$. Since $p \nmid 14$, we have $p|b$. But $\gcd(b,15)=1$ means $p \nmid b$. Contradiction.
    \end{itemize}
    
    Therefore $\gcd(c,210)=1$ $\checkmark$
    
    \textbf{``Only if'' direction ($\Leftarrow$)}: 
    
    If $\gcd(c,210)=1$, then $\gcd(a,14)=1$ AND $\gcd(b,15)=1$.
    
    This is stronger than Statement II (which uses ``or''). Since Statement II is true, and this requires both conditions, we need to verify.
    
    Actually, we can use GCD property:
    $$\gcd(c,14) = \gcd(15a+14b, 14) = \gcd(15a, 14) = \gcd(a,14)$$
    
    (since $\gcd(15,14)=1$)
    
    Similarly: $\gcd(c,15) = \gcd(b,15)$
    
    Since $\gcd(14,15)=1$:
    $$\gcd(c,210) = \gcd(c,14 \cdot 15) = \gcd(c,14) \cdot \gcd(c,15) = \gcd(a,14) \cdot \gcd(b,15)$$
    
    If $\gcd(c,210)=1$, then $\gcd(a,14) \cdot \gcd(b,15) = 1$.
    
    Since both are positive integers, this requires $\gcd(a,14)=1$ AND $\gcd(b,15)=1$ $\checkmark$
    
    Therefore Statement III is \textbf{TRUE}.
    
    \item \textbf{Final answer (5 seconds)}
    
    Statements II and III are true, Statement I is false.
    
    Answer: \textbf{(E) II and III only}
\end{enumerate}

\subsection*{B. Alternative Counterexamples for Statement I}

Other counterexamples that work:
\begin{itemize}
    \item $a=1, b=3$: $c = 15(1) + 14(3) = 15+42 = 57 = 3 \cdot 19$
    \begin{itemize}
        \item $\gcd(1,14)=1$ $\checkmark$, $\gcd(3,15)=3 \neq 1$
        \item $\gcd(57, 210) = 3 \neq 1$ (Statement I violated)
    \end{itemize}
    
    \item $a=1, b=5$: $c = 15(1) + 14(5) = 15+70 = 85 = 5 \cdot 17$
    \begin{itemize}
        \item $\gcd(1,14)=1$ $\checkmark$, $\gcd(5,15)=5 \neq 1$
        \item $\gcd(85, 210) = 5 \neq 1$ (Statement I violated)
    \end{itemize}
    
    \item $a=3, b=5$: $c = 15(3) + 14(5) = 45+70 = 115 = 5 \cdot 23$
    \begin{itemize}
        \item $\gcd(3,14)=1$ $\checkmark$, $\gcd(5,15)=5 \neq 1$
        \item $\gcd(115, 210) = 5 \neq 1$ (Statement I violated)
    \end{itemize}
\end{itemize}

All of these work as counterexamples!

\subsection*{C. Rigorous Proof Using GCD Properties}

The elegant approach uses:
$$\gcd(c,210) = \gcd(a,14) \cdot \gcd(b,15)$$

This immediately proves:
\begin{itemize}
    \item Statement III: $\gcd(c,210)=1 \Leftrightarrow \gcd(a,14)=\gcd(b,15)=1$
    \item Statement II: If $\gcd(c,210)=1$, then $\gcd(a,14)=\gcd(b,15)=1$, which implies at least one is 1
    \item Statement I is false: If only one of $\gcd(a,14)$ or $\gcd(b,15)$ is 1, the product may not be 1
\end{itemize}
\end{solutionbox}

\begin{competitionbox}
\subsection*{A. Time Management}
\begin{itemize}[leftmargin=*]
    \item \textbf{Estimated Time}: 5-7 minutes total
    \begin{itemize}
        \item Quick approach (counterexample + answer choice elimination): 3-4 minutes
        \item Rigorous approach (prove all statements): 7-9 minutes
    \end{itemize}
    \item \textbf{Time Checkpoints}: 
    \begin{itemize}
        \item At 1 minute: Should have converted to $15a+14b=c$
        \item At 2 minutes: Should have found counterexample to Statement I
        \item At 3 minutes: Should have eliminated (A), (B), (C) and focused on II vs III
        \item If stuck: Try simple values like $a=1, 2, 3$ with $b=1, 2, 3$
    \end{itemize}
    \item \textbf{Priority Level}: High - Mid-to-late band problem testing important techniques
\end{itemize}

\subsection*{B. Risk Assessment}
\begin{itemize}[leftmargin=*]
    \item \textbf{Confidence Level}: 85-90\% after finding counterexample and using answer choice logic
    \item \textbf{Common Errors}:
    \begin{itemize}
        \item Misreading ``or'' as ``and'' (or vice versa)
        \item Not checking counterexample arithmetic carefully
        \item Confusing ``if-then'' with ``if-and-only-if''
        \item Spending too long trying to prove Statement I (just find counterexample!)
    \end{itemize}
    \item \textbf{Abandonment Criteria}: If spending more than 8 minutes, make educated guess
    \item \textbf{Flag for Review}: If counterexample arithmetic is uncertain
\end{itemize}

\subsection*{C. Competition Context}
\begin{itemize}[leftmargin=*]
    \item \textbf{Problem 15/25 Strategy}: Transition to late-band. This separates students who can think logically about number theory from those who can't.
    \item \textbf{Point Value}: 6 points (standard AMC12), with 1.5 points for leaving blank
    \item \textbf{Difficulty Assessment}: Mid-to-late band. Tests:
    \begin{itemize}
        \item GCD properties and number theory reasoning
        \item Logical analysis (or, and, if-then, if-and-only-if)
        \item Counterexample construction
        \item Strategic use of answer choices
        \item Modular arithmetic (for rigorous approach)
    \end{itemize}
    \item \textbf{Strategic Value}: This problem type (logical statements about number theory) appears regularly on AMC. Master the technique of testing with simple counterexamples.
\end{itemize}
\end{competitionbox}

\begin{keyinsightbox}
\textbf{Key Insight}: The fundamental relationship is:
$$\gcd(c,210) = \gcd(a,14) \cdot \gcd(b,15)$$

This comes from:
\begin{itemize}
    \item $c = 15a + 14b$
    \item $\gcd(c,14) = \gcd(15a+14b, 14) = \gcd(15a, 14) = \gcd(a,14)$ (using $\gcd(15,14)=1$)
    \item $\gcd(c,15) = \gcd(b,15)$ (similarly)
    \item $\gcd(c,210) = \gcd(c,14 \cdot 15) = \gcd(c,14) \cdot \gcd(c,15)$ (using $\gcd(14,15)=1$)
\end{itemize}

From this relationship:
\begin{itemize}
    \item \textbf{Statement III is true}: $\gcd(c,210)=1 \Leftrightarrow \gcd(a,14) \cdot \gcd(b,15)=1 \Leftrightarrow \gcd(a,14)=\gcd(b,15)=1$
    \item \textbf{Statement II is true}: It's the ``only if'' part of Statement III, which we just proved
    \item \textbf{Statement I is false}: If only ONE of the GCDs is 1, the product may not be 1. Example: $\gcd(a,14)=2, \gcd(b,15)=1$ gives $\gcd(c,210)=2 \neq 1$
\end{itemize}

The strategic insight for competition: Find a simple counterexample to Statement I (like $a=2, b=2$), eliminate three answer choices, then use the answer choice structure to deduce the rest!
\end{keyinsightbox}

\begin{warningbox}
\textbf{Warning}: Common Pitfalls to Avoid

\begin{itemize}
    \item \textbf{Misreading logical connectives}: 
    \begin{itemize}
        \item ``Or'' means AT LEAST ONE (inclusive or), not exactly one
        \item ``And'' means BOTH must be true
        \item ``If and only if'' means BOTH directions: $P \Rightarrow Q$ AND $Q \Rightarrow P$
    \end{itemize}
    
    \item \textbf{Arithmetic errors in counterexamples}: 
    \begin{itemize}
        \item Always verify: $c = 15a + 14b$ computed correctly
        \item Check GCDs carefully: $\gcd(58, 210) = 2$, not 1
        \item Don't assume - calculate!
    \end{itemize}
    
    \item \textbf{Trying to prove false statements}: 
    \begin{itemize}
        \item If Statement I seems hard to prove, try to disprove it!
        \item One counterexample is sufficient to show a statement is not ``necessarily true''
    \end{itemize}
    
    \item \textbf{Confusing Statement II and III}: 
    \begin{itemize}
        \item Statement II: $\gcd(c,210)=1 \Rightarrow \gcd(a,14)=1$ OR $\gcd(b,15)=1$ (or both)
        \item Statement III: $\gcd(c,210)=1 \Leftrightarrow \gcd(a,14)=1$ AND $\gcd(b,15)=1$
        \item III is STRONGER (requires both, not just one)
    \end{itemize}
    
    \item \textbf{Not using answer choice structure}: 
    \begin{itemize}
        \item Once you eliminate (A), (B), (C), notice that III appears in BOTH remaining choices
        \item This means III must be true - no need to fully prove it if pressed for time!
    \end{itemize}
    
    \item \textbf{Time mismanagement}: 
    \begin{itemize}
        \item Don't spend 5 minutes trying to rigorously prove all three statements
        \item Quick strategic approach: 2 minutes to find counterexample, 1 minute to narrow choices
    \end{itemize}
\end{itemize}
\end{warningbox}

\begin{tipbox}
\textbf{Pro Tip}: Competition Strategy

\begin{itemize}
    \item \textbf{Test Statement I first}: It's often the easiest to disprove with a counterexample. This eliminates half the answer choices!
    
    \item \textbf{Use small values}: Try $a=1, 2, 3$ and $b=1, 2, 3$. Don't overthink - simple counterexamples often work.
    
    \item \textbf{Answer choice leverage}: 
    \begin{itemize}
        \item After eliminating (A), (B), (C), notice that (D) and (E) both contain III
        \item Therefore III must be true - you can assume this without proof!
        \item Only need to determine if II is true to choose between (D) and (E)
    \end{itemize}
    
    \item \textbf{Contrapositive is your friend}: Statement II can be proven via contrapositive, which is often easier than direct proof.
    
    \item \textbf{Know GCD properties}: 
    \begin{itemize}
        \item $\gcd(ma+nb, m) = \gcd(nb, m)$ when $m|ma$
        \item When $\gcd(m,n)=1$: $\gcd(k, mn) = \gcd(k,m) \cdot \gcd(k,n)$
    \end{itemize}
    
    \item \textbf{Verify counterexamples}: Always check that your counterexample:
    \begin{itemize}
        \item Satisfies the original equation
        \item Satisfies the condition of the statement
        \item Violates the conclusion of the statement
    \end{itemize}
    
    \item \textbf{Time allocation}: 
    \begin{itemize}
        \item 2 minutes: Find counterexample to Statement I
        \item 1 minute: Use answer choices to narrow to (D) or (E)
        \item 2-3 minutes: Determine if Statement II is true
        \item Total: 5-6 minutes
    \end{itemize}
    
    \item \textbf{Pattern recognition}: When you see GCD relationships in a linear equation, think about modular arithmetic and the GCD product rule.
\end{itemize}
\end{tipbox}

\section*{Final Answer}

The statements that are necessarily true are:
$$\boxed{\textbf{(E) } \text{II and III only}}$$

\textbf{Summary}:
\begin{itemize}
    \item \textbf{Statement I is FALSE}: Counterexample $a=2, b=2$ gives $c=58$ with $\gcd(2,15)=1$ but $\gcd(58,210)=2 \neq 1$
    \item \textbf{Statement II is TRUE}: Can be proven via contrapositive or using $\gcd(c,210)=\gcd(a,14) \cdot \gcd(b,15)$
    \item \textbf{Statement III is TRUE}: $\gcd(c,210)=1 \Leftrightarrow \gcd(a,14)=\gcd(b,15)=1$ (follows from GCD product rule)
\end{itemize}

\section*{Confidence Assessment}
\begin{itemize}[leftmargin=*]
    \item \textbf{Confidence Level}: 90\% - Counterexample verified, answer choice logic is sound, rigorous proof available for Statements II and III
    \item \textbf{Verification Applied}: Level 4-5 (Standard to Thorough) - Checked counterexample arithmetic and logical reasoning
    \item \textbf{Flag for Review}: No - Solution is complete and verified
    \item \textbf{Time Spent}: Estimated 5-6 minutes (optimal for P15)
\end{itemize}

\section*{Summary of Solution Path}

\textbf{Strategic Approach (5-6 minutes):}
\begin{enumerate}
    \item Convert to $15a + 14b = c$
    \item Test Statement I: Find counterexample $a=2, b=2 \Rightarrow c=58$
    \item Verify: $\gcd(2,15)=1$ $\checkmark$, but $\gcd(58,210)=2 \neq 1$ (Statement I false)
    \item Eliminate answer choices (A), (B), (C)
    \item Observe: Statement III appears in both (D) and (E), so it must be true
    \item Determine Statement II: Use contrapositive or GCD properties to show it's true
    \item Answer: (E) II and III only
\end{enumerate}

\textbf{Rigorous Approach (7-9 minutes):}
\begin{enumerate}
    \item Use GCD property: $\gcd(c,210) = \gcd(a,14) \cdot \gcd(b,15)$
    \item Prove directly:
    \begin{itemize}
        \item Statement I false: Product of two GCDs (one being >1) may not be 1
        \item Statement II true: Follows from Statement III's ``only if'' part
        \item Statement III true: Product equals 1 $\Leftrightarrow$ both factors equal 1
    \end{itemize}
    \item Answer: (E) II and III only
\end{enumerate}

\end{document}

